\documentclass[11pt,a4paper]{article}
% Latin Modern (T1) when installed; otherwise vector Computer Modern (OT1), which matches the figures.
\IfFileExists{lmodern.sty}{\usepackage[T1]{fontenc}\usepackage{lmodern}}{}
\usepackage[margin=2.4cm]{geometry}
\usepackage{amsmath,amssymb,amsthm,bm}
\usepackage{booktabs}
\usepackage{graphicx}
\usepackage{adjustbox}
\usepackage{caption}
\usepackage{float}
\usepackage{xcolor}
\usepackage{enumitem}
\usepackage{microtype}
\usepackage[colorlinks=true,linkcolor=blue!50!black,citecolor=blue!50!black,urlcolor=blue!50!black]{hyperref}
\usepackage[capitalise,noabbrev]{cleveref}

\captionsetup{font=small,labelfont=bf,margin=0.5em}
\setlist{itemsep=2pt,topsep=3pt}
\graphicspath{{figures/}}
\newtheorem{proposition}{Proposition}
\newtheorem{remark}{Remark}

\newcommand{\R}{\mathbb{R}}
\newcommand{\X}{\mathcal{X}}
\newcommand{\B}{\mathcal{B}}
\newcommand{\U}{\mathcal{U}}
\newcommand{\Sset}{\mathcal{S}}
\newcommand{\Fset}{\mathcal{F}}
\newcommand{\hc}{h_{\mathrm{cone}}}
\newcommand{\zd}{z_{\mathrm{des}}}
\newcommand{\tabfit}[1]{\begin{adjustbox}{max width=\linewidth}\input{tables/#1}\end{adjustbox}}

\input{numbers.tex}

\title{\textbf{PS2-RL Phase I for Approach-Cone Landing}\\[4pt]
\large Base-set certification, safe-arrival policy training,\\ and closed-loop verification of trajectories and actuator commands}
\author{Jungwon Kim\\ \small Seoul National University}
\date{30 September 2026}

\begin{document}
\maketitle

\begin{abstract}
\noindent
This report documents Phase~I of PS2-RL for the approach-cone landing task: the construction and
certification of a base set and base controller inside a smooth approach cone, and the training and
verification of a safe-arrival policy that steers the quadrotor into that base set. The base set is a
9-D LQR ellipsoid about a hover point $\zd=\zDes$\,m above the pad, at level $c_{\B}=\cB$. Among
the level bounds, the numerically certified one-step Lyapunov decrease ($c_{\mathrm{Lyap}}=\cLyap$) binds;
input feasibility gives $c_{\U}=\cU$ and cone containment $\cConeZ$. Six safe-arrival policies (TD3 and SAC,
three seeds each, 5M steps) raise the weighted recoverability of held-out initial states from
$\muUntrained$ (untrained) and $\muLQR$ (base controller used everywhere) to \muTDthree{} (TD3) and
\muSAC{} (SAC). The selected policy (\runChosen) reaches $\hat\mu_w=\muChosen$, which is \fracOfCeil\% of a
viability ceiling of $\muCeil$ estimated by trajectory optimisation: almost every remaining failure is a
state from which no admissible input sequence can stay inside the cone. Closed-loop case studies confirm
the intended behaviour (climb-and-turn-in recoveries, hand-off to the base controller, forward invariance
afterwards) and expose two issues that matter for deployment: the model has no ground plane, and the
policy banks hard (median peak tilt $\approx\tiltMedGen^\circ$, 95th percentile up to $\tiltPninetyfiveGen^\circ$).
\end{abstract}

\tableofcontents
\newpage

%======================================================================
\section{Scope and summary}
%======================================================================

PS2-RL \cite{ps2rl} trains a task policy through a differentiable control-invariant layer (CIL) whose
safety certificate is a backup control barrier function. The backup policy $\pi_b$ combines a
certified base controller $\pi_{\B}$ on a small base set $\B$ with a learned \emph{safe-arrival}
policy $\pi_{\mathrm{SA}}$ outside it. Phase~I trains $\pi_{\mathrm{SA}}$; Phase~II trains the task
policy through the CIL. Moving from the paper's powerloop task to landing changes the safe set from a
ceiling to an approach cone. Because a cone is bounded horizontally, the base set must also regulate
horizontal position, which in turn changes the base controller, its certificate and the Phase~I design region.
The treatment follows the Phase~I reformulation note for the landing task \cite{landingnote}.

\paragraph{Summary of findings.}
\begin{itemize}
\item \textbf{Certified base set.} The yaw chain of the powerloop LQR weights limits the admissible level to
$c\le\cChartpl$ via the quaternion chart. Re-weighting only the (decoupled) yaw chain lifts every bound
above $15$; the binding one is the Lyapunov decrease, $c_{\mathrm{Lyap}}=\cLyap$. We use $c_{\B}=\cB$
($\cBratio\,c_{\mathrm{Lyap}}$), and $\zd=\zDes$\,m is the lowest hover height at which neither the cone
nor the ground clearance binds.
\item \textbf{Training.} Both TD3 and SAC converge in every seed; runs take $\approx\wallMin$\,min each on a
laptop GPU. At the full budget TD3 is marginally but consistently better (\muTDthree{} vs.\ \muSAC{}); SAC
reaches its plateau with fewer samples.
\item \textbf{The learned policy works, up to the physics.} On held-out test states the selected policy
arrives safely from $\mGenChosen$ of the general region, $\mEdgeChosen$ of the low-altitude cone edge and
all of the capture shell, against $\mGenLQR$ / $\mEdgeLQR$ / 1 for the base controller alone. There are no
timeouts: every failure is a crash. Of the \edgeFailChosen{} edge failures, \edgeShared{} are shared by all six
policies, and trajectory optimisation finds no safe input sequence for \edgeNonviable{} of those.
\item \textbf{Open issues.} Some successful recoveries (0.2--0.8\%) dip a few centimetres below the pad
plane because the model has no ground. The policy also uses large attitudes, which the SITL rate loop may not
reproduce (\cref{sec:limitations}).
\end{itemize}

%======================================================================
\section{System model}
%======================================================================

\subsection{Continuous-time dynamics}
The state is $x=(p,v,q)\in\R^3\times\R^3\times\mathbb{S}^3$ (position, inertial velocity, unit quaternion
$q=(q_w,q_x,q_y,q_z)$, body-to-world) and the input is $u=(a,\omega)$, a mass-normalised collective thrust and
body rates:
\begin{equation}
\dot p = v,\qquad \dot v = -g e_3 + a\,R(q)e_3,\qquad \dot q = \tfrac12\,\Xi(q)\,\omega,\qquad
\Xi(q)=\begin{bmatrix}-q_x&-q_y&-q_z\\ q_w&-q_z&q_y\\ q_z&q_w&-q_x\\ -q_y&q_x&q_w\end{bmatrix}.
\label{eq:dyn}
\end{equation}
Body rates are commanded directly and tracked instantaneously; the inner rate loop is not modelled.
The input set is the box $\U=[0,a_{\max}]\times[-\bar\omega,\bar\omega]^3$ with $a_{\max}=4g$ and
$\bar\omega=18$\,rad/s (\cref{tab:params}). Pad-relative coordinates are
$\Delta p=p_{xy}-p_{\mathrm{pad},xy}$ and $\zeta=z-z_{\mathrm{pad}}$.

\subsection{Discretisation}
Training, evaluation, certification and the CIL all use the same sampled model: an explicit Euler step
followed by quaternion re-normalisation,
\begin{equation}
F(x,u)=\mathcal{N}\big(x+\Delta\,f(x,\mathrm{sat}_{\U}(u))\big),\qquad \Delta=0.02\,\mathrm{s},
\label{eq:euler}
\end{equation}
where $f$ is the right-hand side of \eqref{eq:dyn} and $\mathcal{N}$ normalises the quaternion block. The
backup horizon is $N=T/\Delta=100$ steps ($T=2$\,s). For a policy $\pi$ we write
$\Phi^{\pi}_{k+1}(x)=F(\Phi^{\pi}_k(x),\pi(\Phi^{\pi}_k(x)))$, $\Phi^\pi_0(x)=x$. All statements below concern
this sampled system.

\begin{table}[t]
\centering\small
\caption{Model and task parameters.}
\label{tab:params}
\input{tables/params}
\end{table}

%======================================================================
\section{Safe set: a smooth one-sided cone}
%======================================================================

The safe set is the approach cone
\begin{equation}
\hc(x)=r_0+\tan\theta\,\zeta-\sqrt{\|\Delta p\|^2+\varepsilon^2},\qquad
\Sset=\{x:\hc(x)\ge 0\},\qquad \Fset=\X\setminus\Sset,
\label{eq:cone}
\end{equation}
with $r_0=0.5$\,m, $\theta=30^\circ$ and $\varepsilon=0.05\,r_0$. Three properties are used throughout:
\begin{enumerate}
\item $\hc\in C^\infty$ with $\nabla_p\hc=(-\Delta p/\sqrt{\|\Delta p\|^2+\varepsilon^2},\ \tan\theta)$. The
$\varepsilon$-smoothing is not cosmetic: the unsmoothed cone is non-differentiable on its axis, which is
exactly where the base-set equilibrium lies, and the CIL differentiates $\hc$ along backup rollouts.
\item $\hc$ is globally Lipschitz with constant $\sec\theta$, and concave in position (affine minus a
convex function). The concavity is used to certify cone containment exactly (\cref{sec:levels}).
\item The square-root form is one-sided: $\hc\ge0\Rightarrow\zeta\ge-(r_0-\varepsilon)/\tan\theta$. The
squared form $(r_0+\tan\theta\,\zeta)^2-\|\Delta p\|^2$ would describe a double cone with a spurious safe
region below the apex.
\end{enumerate}
The cross-section at altitude $\zeta$ is the disk $\|\Delta p\|\le R_\varepsilon(\zeta)=\sqrt{(r_0+\tan\theta\,\zeta)^2-\varepsilon^2}$.
There is no ground constraint; the consequences are discussed in \cref{sec:limitations}.

%======================================================================
\section{Base set and base controller}
%======================================================================

\subsection{Equilibrium and error coordinates}
A set that is unbounded in $(x,y)$ cannot lie inside the bounded cross-sections of the cone, so the
powerloop's 7-D hover error, which omits horizontal position, is not admissible. We take hover above the
pad as the equilibrium,
\begin{equation}
x^\star=(p_{\mathrm{pad}}+\zd e_3,\ 0,\ q_I),\qquad u^\star=(g,0,0,0),
\end{equation}
which lies in the interior of $\Sset$ ($\hc(x^\star)=r_0+\tan\theta\,\zd-\varepsilon>0$) and of $\U$
($0<g<a_{\max}$). On $\{q_w\neq0\}$ the 9-D error coordinates are
\begin{equation}
e(x)=\big(\Delta p_x,\ \Delta p_y,\ \zeta-\zd,\ v_x,\ v_y,\ v_z,\ \phi\big)\in\R^9,\qquad
\phi=2\,\mathrm{sgn}(\bar q_w)\,\bar q_{xyz},\ \ \bar q=q^{-1},
\label{eq:err}
\end{equation}
so $\phi\approx-\vartheta$ at small angles, where $\vartheta$ is the rotation vector of $q$.

\subsection{Linearisation and controllability}
At small angles $R(q)e_3\approx e_3+\vartheta\times e_3=(-\phi_y,\ \phi_x,\ 1)^\top$, and the body rates
coincide with $\dot\vartheta=-\dot\phi$. With $\delta a=a-g$, the linearisation of \eqref{eq:dyn} about
$(x^\star,u^\star)$ in the coordinates \eqref{eq:err} is
\begin{equation}
\dot e = A e + B\,\delta u,\qquad
A=\begin{bmatrix}0_3 & I_3 & 0_3\\ 0_3 & 0_3 & G\\ 0_3 & 0_3 & 0_3\end{bmatrix},\quad
G=\begin{bmatrix}0&-g&0\\ g&0&0\\ 0&0&0\end{bmatrix},\quad
B=\begin{bmatrix}0_{3\times1} & 0_3\\ e_3 & 0_3\\ 0_{3\times1} & -I_3\end{bmatrix},
\label{eq:lin}
\end{equation}
with $\delta u=(\delta a,\omega)$. The state splits into four integrator chains, each ending in a distinct input:
\[
\Delta p_x\to v_x\to\phi_y\leftarrow\omega_y,\qquad \Delta p_y\to v_y\to\phi_x\leftarrow\omega_x,\qquad
\zeta\to v_z\leftarrow\delta a,\qquad \phi_z\leftarrow\omega_z,
\]
of lengths $3+3+2+1=9$. The pair $(A,B)$ is therefore in Brunovsky form and controllable; the discrete
controllability matrix has full rank 9 numerically as well (\cref{tab:extents}).

\subsection{Discretisation of the linear model and LQR design}
The linear model is discretised with the same explicit Euler rule as the plant,
\begin{equation}
A_d=I_9+\Delta A,\qquad B_d=\Delta B .
\end{equation}
This is not an approximation of convenience: because \eqref{eq:euler} is an Euler step, its Jacobian at
$(x^\star,u^\star)$ in the coordinates \eqref{eq:err} is exactly $(A_d,B_d)$. The re-normalisation
$\mathcal{N}$ only acts at second order, since $\dot q$ is tangent to $\mathbb{S}^3$ at $q_I$. The LQR
therefore stabilises the true linearisation of the sampled plant, with no zero-order-hold mismatch.
With diagonal weights $Q\succ0$ and $R\succ0$, $P\succ0$ solves the discrete algebraic Riccati equation
\begin{equation}
P=A_d^\top PA_d-A_d^\top PB_d\,(R+B_d^\top PB_d)^{-1}B_d^\top PA_d+Q,\qquad
K=(R+B_d^\top PB_d)^{-1}B_d^\top PA_d,
\end{equation}
and $A_d-B_dK$ is Schur, with spectral radius $\rho=\rhoCL$. The base controller and base set are
\begin{equation}
\begin{gathered}
\pi_{\B}(x)=\mathrm{sat}_{\U}\big(u^\star-Ke(x)\big),\qquad V(x)=e(x)^\top Pe(x),\\
\B_c=\{x:V(x)\le c\},\qquad h_{\B}(x)=c-V(x).
\end{gathered}
\label{eq:base}
\end{equation}
The gain is listed in \cref{tab:gain} (Appendix).

\paragraph{Choice of weights.}
The position and velocity rows reuse the powerloop weights, with $q_x=q_y=q_\zeta=1$
(\cref{tab:lqrw}). The four chains decouple in $P$: numerically $|P_{xy}|,|P_{x\zeta}|<10^{-13}$. The yaw
chain can therefore be tuned without touching the others. With the powerloop yaw weights
($q_{\phi_z}=0.16$, $r_{\omega_z}=0.004$), the ellipsoid admits yaw errors beyond what the quaternion chart
can represent: $c_{\mathrm{chart}}=\cChartpl$, and the one-step Lyapunov decrease fails near $c\approx6.7$.
Raising the yaw cost to $q_{\phi_z}=0.5$ and $r_{\omega_z}=0.04$ (a slightly gentler yaw correction) moves
both bounds above 15 and leaves every other row of $K$ unchanged.

\begin{table}[t]
\centering\small
\begin{minipage}[t]{0.44\linewidth}\centering
\caption{LQR weights. $^\dagger$\,changed.}
\label{tab:lqrw}
\input{tables/lqr_weights}
\end{minipage}\hfill
\begin{minipage}[t]{0.52\linewidth}\centering
\caption{Base set at $c_{\B}=\cB$ and closed-loop facts.}
\label{tab:extents}
\input{tables/extents}
\end{minipage}
\end{table}

\subsection{Certifying the level \texorpdfstring{$c$}{c}}\label{sec:levels}
For an ellipsoid $\{e^\top Pe\le c\}$, a vector $w$ and a matrix $E$,
\begin{equation}
\sup_{e^\top Pe\le c} w^\top e=\sqrt{c\,w^\top P^{-1}w},\qquad
\sup_{e^\top Pe\le c}\|Ee\|=\sqrt{c\,\lambda_{\max}(EP^{-1}E^\top)} .
\label{eq:support}
\end{equation}
Every support value scales with $\sqrt c$, so each bound below is the largest level satisfying a
monotone condition.
\begin{itemize}
\item \textbf{Input feasibility.} $c_{\U}=\min_i \bar u_i^2/(K_iP^{-1}K_i^\top)$ with
$\bar u_i=\min(u^\star_i-u_{\min,i},\,u_{\max,i}-u^\star_i)$. For $c\le c_{\U}$ the saturation in
\eqref{eq:base} is inactive on $\B_c$, so $\pi_{\B}$ is smooth there.
\item \textbf{Chart validity.} $c_{\mathrm{chart}}=4/\lambda_{\max}(E_\phi P^{-1}E_\phi^\top)$ keeps
$\|\phi\|<2$, so $q_w\neq0$ and \eqref{eq:err} remains a valid chart.
\item \textbf{Cone containment (exact).} Because $\hc$ is concave in position, its minimum over the
(convex) position projection of $\B_c$ lies on the boundary. We parametrise that boundary as
$e_p=\sqrt c\,Lu$ with $LL^\top=E_pP^{-1}E_p^\top$ and $u\in\mathbb{S}^2$, evaluate $\hc$ on a
$360\times720$ grid, and bisect on $c$. The corner bound of the note, which takes the radial and altitude
extremes simultaneously, is valid but conservative ($15.20$ vs.\ $\cConeZ$ at $\zd=\zDes$\,m;
\cref{tab:zsweep}).
\item \textbf{Ground clearance.} $c_{\mathrm{ground}}=(\zd-z_{\mathrm{clear}})^2/(P^{-1})_{33}$ keeps $\B_c$
at least $z_{\mathrm{clear}}=0.5$\,m above the pad. This bound is not in the note; it keeps the hover set
out of ground effect.
\item \textbf{Lyapunov decrease.} $c_{\mathrm{Lyap}}$ is the largest $c$ such that
$V(F(x,\pi_{\B}(x)))<V(x)$ on every level set $\{V=c'\}$, $c'\le c$, for the \emph{nonlinear} step
\eqref{eq:euler}. Existence follows from the Schur property (discrete Lyapunov indirect method); its
value is estimated numerically. Random directions in 9-D miss the worst case (large tilt combined with
falling), so we run projected gradient ascent of $V(x^+)/V(x)$ over each level set from 4096 random
starts and bisect on $c$ (\cref{fig:bounds}b). This is a numerical check, not a proof.
\end{itemize}

\begin{proposition}[Base-pair certificate, landing note Prop.~1]
Let
\[\bar c=\min\{c_{\U},\,c_{\mathrm{chart}},\,c_{\mathrm{cone}},\,c_{\mathrm{ground}},\,c_{\mathrm{Lyap}}\}.\]
For every $c\le\bar c$:
\begin{enumerate}[label=(\roman*),nosep]
\item $\B_c\subseteq\Sset$;
\item $\pi_{\B}(\B_c)\subseteq\U$, without saturation;
\item $\B_c$ is forward invariant under $x\mapsto F(x,\pi_{\B}(x))$.
\end{enumerate}
\end{proposition}

\begin{table}[t]
\centering\small
\caption{Level bounds at $\zd=\zDes$\,m, for the landing weights and, for comparison, with the powerloop yaw
weights. The Lyapunov value for the powerloop yaw weights comes from a random-direction sweep.}
\label{tab:bounds}
\tabfit{bounds}
\end{table}

\begin{figure}[t]
\centering
\includegraphics[width=\linewidth]{fig_level_bounds.pdf}
\caption{\textbf{Choosing $\zd$ and $c_{\B}$.} (a)~Only the cone and ground bounds depend on the hover
height. They rise with $\zd$ and stop binding at $\zd\approx\zSweet$\,m, above which
$\bar c=c_{\mathrm{Lyap}}$. The corner bound (dashed) would have suggested a higher hover point.
(b)~Worst one-step ratio $V(x^+)/V(x)$ found by adversarial search on each level set. It crosses 1 at
$c_{\mathrm{Lyap}}=\cLyap$; the chosen $c_{\B}=\cB$ keeps a margin of $\approx3\%$ in the ratio.}
\label{fig:bounds}
\end{figure}

\paragraph{Choice of $\zd$.} The hover linearisation is translation-invariant, so $P$, $K$, $c_{\U}$,
$c_{\mathrm{chart}}$ and $c_{\mathrm{Lyap}}$ do not depend on $\zd$; only $c_{\mathrm{cone}}(\zd)$ and
$c_{\mathrm{ground}}(\zd)$ do, and both increase with $\zd$ (\cref{fig:bounds}a, \cref{tab:zsweep}). The
admissible level $\bar c(\zd)$ therefore grows until the $\zd$-independent bound takes over, at
$\zd\approx\zSweet$\,m. Beyond that point a higher hover set is no larger; it only lengthens the climb from
the low-altitude states that decide recoverability. We use $\zd=\zDes$\,m, just above the knee. Ground
clearance binds before the cone does for any cone at least as wide as $r_0=0.5$\,m, $\theta=20^\circ$.

\paragraph{Choice of $c_{\B}$.} $c_{\B}=\cB=\cBratio\,c_{\mathrm{Lyap}}$ leaves margin for the numerical
nature of $c_{\mathrm{Lyap}}$; the paper's powerloop used $8/12.38=0.65$ of its bound. At this level the base
set reaches $\pm\extXY$\,m horizontally, $\pm\extZ$\,m vertically and $\pm\extV$\,m/s in horizontal velocity,
and its lowest point is $\lowestB$\,m above the pad (\cref{tab:extents}). The large roll/pitch extent
($\max|\phi_x|=1.44$) is attained only jointly with position and velocity errors along the ellipsoid's
principal directions.

\paragraph{Numerical verification.} $2\times10^5$ states sampled uniformly in $\B$ all satisfy
$\hc\ge\containMinH$\,m. $4096$ rollouts of $\pi_{\B}$ for $500$ steps from the boundary $\{V=c_{\B}\}$, including
the adversarial worst direction, never exceed $V=\invMaxV<c_{\B}$ and keep $\hc\ge\invMinH$\,m.

%======================================================================
\section{The Phase I problem}
%======================================================================

\subsection{Composite backup and the backup-induced set}
The backup policy and the set it certifies are
\begin{equation}
\pi_b(x)=\begin{cases}\pi_{\B}(x),&x\in\B\\ \pi_{\mathrm{SA}}(x),&x\notin\B\end{cases},\qquad
\mathcal{C}_N(\pi_b)=\{x:\Phi^{\pi_b}_k(x)\in\Sset\ (0\le k\le N),\ \Phi^{\pi_b}_N(x)\in\B\}.
\label{eq:CN}
\end{equation}
With first hitting times $\tau_{\B}$ and $\tau_{\Fset}$, and $(\B,\pi_{\B})$ certified,
$\mathcal{C}_N(\pi_b)=\{x:\tau_{\B}(x)\le N,\ \tau_{\B}(x)<\tau_{\Fset}(x)\}$ (landing note, Lemma~4). Once
the trajectory reaches $\B$ it stays there and inside $\Sset$. Membership in $\mathcal{C}_N$ is thus a
first-hit event, which is what the value function below measures and what all evaluations in this report
count. Phase~I solves
\begin{equation}
\max_{\pi_{\mathrm{SA}}}\ \mu\big(\mathcal{C}_N(\pi_b)\cap(\Omega\setminus\B)\big),
\label{eq:sp2}
\end{equation}
where $\Omega$ is a design region and $\mu$ a reference measure on it (\cref{sec:design}).

\subsection{Safe-arrival value function}
With indicators $b=\mathbf{1}_{\B}$, $f=\mathbf{1}\{\hc<0\}$ and $c=1-b-f$, the safe-arrival Q-function and its
Bellman equation are
\begin{equation}
Q^\pi_\beta(x,u)=\sum_{k\ge0}\beta^k\Big(\prod_{i<k}c(x_i)\Big)b(x_k),\qquad
Q^\pi_\beta(x,u)=b(x)+\beta\,c(x)\,Q^\pi_\beta\big(F(x,u),\pi(F(x,u))\big).
\label{eq:value}
\end{equation}
The optimal value is $V^\star_\beta(x)=\beta^{d(x)}$, where $d(x)$ is the minimum number of steps to
reach $\B$ safely. A discount-optimal policy decreases $d$ by one per step, so it maximises
\eqref{eq:sp2} simultaneously for every horizon $N$ (landing note, Prop.~3). The critic is trained on the target
\begin{equation}
y=\beta\Big(b(x')+c(x')\min_{i=1,2}Q'_i\big(x',\mu'(x')\big)\Big),\qquad x'=F(x,u),
\label{eq:target}
\end{equation}
with twin target critics $Q'_i$ and the target actor's mean action $\mu'$. The critic outputs a sigmoid, so
values stay in $[0,1]$.

\paragraph{Discount factor.} States that can arrive within the horizon ($d\le N$) should keep a value of at
least $v_{\min}$ above failure: $\beta^N\ge v_{\min}\iff\beta\ge v_{\min}^{1/N}$. With $N=100$ and
$v_{\min}=0.3$ this requires $\beta\ge0.988$. We use $\beta=0.99$ ($\beta^{100}=0.366$). The unicycle value
$\beta=0.92$ would give $\beta^{100}\approx2\times10^{-4}$, and the critic could not tell a late arrival
from a crash.

\subsection{Design region, reference measure and sampling}\label{sec:design}
The design region is
$\Omega=\{x\in\Sset:\ 0\le\zeta\le\bar\zeta,\ \|v\|\le\bar v,\ \mathrm{tilt}(q)\le\bar\vartheta,\ |\mathrm{yaw}|\le45^\circ\}$
(\cref{tab:region}). It is covered by three analytically sampled regions (\cref{fig:geometry}); every
sample is rejection-sampled to $\Omega\setminus\B$.
\begin{itemize}
\item \textbf{General (altitude-uniform).} $\zeta\sim\mathcal{U}[0,\bar\zeta]$; then $\Delta p$ uniform by
area on the cross-section disk of radius $R_\varepsilon(\zeta)$; velocity uniform in the ball; body
$z$-axis uniform on the tilt cap; yaw uniform. Sampling uniformly in \emph{volume} would put most samples
high up, where the cone is wide and arrival is easy. Uniform per altitude slice gives each altitude equal mass.
\item \textbf{Edge band.} $0\le\hc\le\eta$, $\zeta\le\zeta_{\mathrm{low}}$, with the radial velocity turned
outward with probability $0.7$. These low-altitude states near the wall are the ones that decide
landing safety (the analogue of the powerloop's near-ceiling states).
\item \textbf{Capture shell.} $c_{\B}<V\le\kappa c_{\B}$: level uniform, direction uniform in the $P$-metric.
\end{itemize}
Training draws from a curriculum mixture that shifts weight from general and shell towards the edge band as
the curriculum scale $s$ grows. The scale increases by 0.1 whenever the arrival rate within $N$ over the
last 200 episodes exceeds 0.8. The general region stays in the mixture at every scale, since optimality is
needed on all of $\mathcal{X}\setminus(\B\cup\Fset)$ and rollouts leave $\Omega$. The reference measure
$\mu=\sum_r w_r\mu_r$ uses weights $1/3/0.5$ (general/edge/shell).

\begin{figure}[t]
\centering
\includegraphics[width=0.92\linewidth]{fig_geometry.pdf}
\caption{\textbf{Geometry of the landing Phase I problem}, as a vertical slice through the cone axis.
Held-out test states of the three regions are shown at their signed radial distance. The base set $\B$
is drawn as its position projection; actual membership also depends on velocity and attitude.}
\label{fig:geometry}
\end{figure}

\begin{table}[t]
\centering\small
\caption{Design region, curriculum mixture and reference measure.}
\label{tab:region}
\input{tables/region}
\end{table}

\subsection{Learners}
Both learners share the critic of \eqref{eq:target}: twin networks with sigmoid output, a Huber loss and
delayed actor and target updates. Outside $\B$ the actor output is used as is. Inside $\B$ the base controller
acts, so actor-loss terms from states in $\B$ are masked (hand-off masking). Actions are squashed into
$\U$ with $a=u_{\mathrm{mid}}+u_{\mathrm{half}}\odot\tanh(\cdot)$. A small action-magnitude penalty
($0.05\,\|\tilde u\|^2$ on the normalised action $\tilde u$) regularises the actor. The actor uses ELU
activations, so that $\pi_b$ is $C^1$ where the CIL integrates its sensitivity.

\paragraph{TD3.} Deterministic actor $\mu_\psi$, trained by maximising $\min_iQ_i(x,\mu_\psi(x))$, with
Gaussian exploration noise of fixed scale during collection.

\paragraph{SAC variant.} A tanh-Gaussian actor is trained with
$L_\pi=\mathbb{E}\big[\alpha\log\pi(\tilde a\,|\,x)-\min_iQ_i(x,a)\big]$, with $\alpha$ auto-tuned towards a
target entropy of $-4$ on the normalised action. The entropy term appears \emph{only} in the actor loss;
the critic keeps the hard backup \eqref{eq:target} on the deterministic mean. A soft backup would add
$-\alpha\log\pi$ as a per-step reward and move the fixed point away from $\beta^{d(x)}$. With positive
entropy, a state that never arrives accrues $\approx\alpha\mathcal{H}/(1-\beta)$ and can outrank a late
arrival. With negative entropy, failure (value 0, episode over) can outrank a long arrival. Either way the
ordering ``arrive sooner $>$ arrive later $>$ never arrive'' that \eqref{eq:sp2} relies on is lost. Keeping
the backup on the mean also makes $Q$ estimate the value of exactly the policy that $\pi_b$ deploys. What
SAC changes is exploration, which becomes state-dependent.

\begin{table}[t]
\centering\small
\caption{Training hyperparameters (shared by both learners unless stated).}
\label{tab:hparams}
\tabfit{hparams}
\end{table}

\subsection{Evaluation protocol}
Two fixed, i.i.d.\ held-out sets are drawn from the three regions at full curriculum scale (validation
and test; 1024/1024/512 states per region). For each region, the recoverability
\begin{equation}
\hat m_r=\frac1{M_r}\sum_{i=1}^{M_r}\mathbf{1}\big[x^{(i)}\in\mathcal{C}_N(\pi_b)\big],\qquad
\hat\mu_w=\frac{\sum_r w_r\hat m_r}{\sum_r w_r},
\end{equation}
is computed by rolling out the deterministic backup policy with first-hit semantics. By Hoeffding's
inequality $|\hat m_r-m_r|\le\sqrt{\ln(2/\delta)/(2M_r)}$ with probability $1-\delta$; this is $\pm\hwEdge$
for $M_r=1024$ at $\delta=0.05$. Checkpoints are selected on validation every 50k steps; all numbers
reported below are on the test split at the selected checkpoint.

%======================================================================
\section{Results}
%======================================================================

\subsection{Training}
\Cref{fig:training} shows the validation recoverability over training. Every run converges.
\begin{itemize}
\item \textbf{SAC} climbs faster: $\hat\mu_w>0.8$ within $\approx0.5$M steps, while TD3 oscillates until
$\approx1.3$M steps.
\item \textbf{Plateau.} Both are flat from $\approx1.5$M steps, so a 2.5M-step budget would suffice.
\item \textbf{Late regression.} SAC regresses slightly late in training (best minus final validation
$\hat\mu_w\approx0.010$, vs.\ $0.004$ for TD3). Its temperature settles near $2\times10^{-4}$
(\cref{fig:trainingdiag}c).
\item \textbf{Curriculum.} The curriculum reaches $s=1$ in every run (\cref{fig:training}f); the crash
rate and the training-time arrival rate plateau together (\cref{fig:trainingdiag}a).
\end{itemize}

\begin{figure}[H]
\centering
\includegraphics[width=\linewidth]{fig_training.pdf}
\caption{\textbf{Validation recoverability during training.} Line: mean over three seeds; band:
min--max across seeds. Step 0 is the untrained actor. End labels give the final mean.}
\label{fig:training}
\end{figure}

\begin{figure}[H]
\centering
\includegraphics[width=\linewidth]{fig_training_diag.pdf}
\caption{\textbf{Training diagnostics}, per run and moving-averaged. (a)~Arrival rate within $N$ of
training episodes, drawn from the curriculum distribution with exploration. (b)~Critic loss. (c)~SAC
temperature.}
\label{fig:trainingdiag}
\end{figure}

\subsection{Recoverability}
\Cref{tab:runs} lists every run and \cref{tab:summary} compares them with baselines.
\begin{itemize}
\item \textbf{Backbones.} Every TD3 seed beats every SAC seed, but only by $0.004$ in $\hat\mu_w$, which
is smaller than the per-region confidence half-width.
\item \textbf{Base-controller baseline.} Using $\pi_{\B}$ everywhere, i.e.\ no safe-arrival policy, is a
strong baseline ($\hat\mu_w=\muLQR$). The learned policy adds most in the edge band, raising it from
$\mEdgeLQR$ to $\mEdgeChosen$.
\item \textbf{Checkpoint choice.} We select \runChosen, which has the highest validation score (0.901).
Its detailed results are in \cref{tab:chosen}. Arrival is fast (mean $\tMeanGen$\,s general,
$\tMeanEdge$\,s edge), so the 2\,s horizon is not what limits recoverability: there are no timeouts.
\end{itemize}

\begin{table}[H]
\centering\small
\caption{Per-run results on the test split at the validation-selected checkpoint. ``best step'' is the
selected checkpoint; ``min'' is the wall time of the 5M-step run on an RTX~4050 laptop GPU (two runs in
parallel).}
\label{tab:runs}
\tabfit{runs}
\end{table}

\begin{table}[H]
\centering\small
\caption{Recoverability of backup policies (test split unless noted; mean $\pm$ std over seeds). The
oracle takes, for each state, the best of the six learned policies. The ceiling counts as unrecoverable only
the shared failures for which trajectory optimisation finds no safe input sequence
(\cref{sec:failures}).}
\label{tab:summary}
\tabfit{summary}
\end{table}

\begin{table}[H]
\centering\small
\caption{Selected policy (\runChosen), test split. $\tau_{\B}$: arrival time into $\B$ of successful
states. Last column: successful arrivals whose path dips below the pad plane $\zeta=0$.}
\label{tab:chosen}
\tabfit{chosen}
\end{table}

\paragraph{Slices of $\mathcal{C}_N$.} \Cref{fig:slices} evaluates the backup on dense grids over the
$(\|\Delta p\|,\zeta)$ half-plane at level attitude, for three velocity conditions.
\begin{itemize}
\item \textbf{At rest.} The learned and base-controller backups are nearly identical (\sliceRest\% vs.\
\sliceRestLQR\% saved).
\item \textbf{With outward velocity.} The base controller loses a wide band along the cone wall that
widens with altitude (\sliceOutLQR\% and \sliceOutDownLQR\%). The learned policy shrinks the lost region to
a sliver next to the wall (\sliceOut\% and \sliceOutDown\%).
\item \textbf{Arrival time} grows smoothly with distance from $\B$; nowhere does it approach the 2\,s horizon.
\end{itemize}

\begin{figure}[H]
\centering
\includegraphics[width=\linewidth]{fig_slices.pdf}
\caption{\textbf{Slices of the backup-induced set $\mathcal{C}_N(\pi_b)$.} Level attitude, yaw 0, radial
velocity as titled. Blue shades: arrival time into $\B$; red: the backup leaves the cone; grey: the initial
state is already in $\B$. Top: learned backup (\runChosen); bottom: base controller alone.}
\label{fig:slices}
\end{figure}

\subsection{Where the failures come from}\label{sec:failures}
The failures are almost identical across policies (\cref{tab:failures}): \edgeShared{} of the
\edgeFailChosen{} edge failures of the selected policy are also failures of all five other policies. To test
whether any controller could save these states, we solve, for each shared failure, an open-loop trajectory
optimisation over the full input box: maximise $\min_k\hc(x_k)$ over the 2\,s horizon, using 9 initialisations
$\times$ 600 Adam steps on a tanh-parametrised input sequence. Only \edgeSharedViable{} of the \edgeShared{} shared edge
failures (and \genSharedViable{} of \genShared{} general ones) can even be kept inside the cone. As a control, the same
procedure succeeds for \ctrlViable{} edge states that all policies save.

The non-viable states share a signature (\cref{fig:viability}): a few centimetres from the wall with an
outward radial velocity. The wall is crossed within $\tfMedian$\,s in the median case (at most $\tfMax$\,s),
faster than thrust and attitude can redirect the vehicle. Counting only these as unrecoverable gives the
ceiling in \cref{tab:summary}. The selected policy reaches \fracOfCeil\% of it, and the per-state oracle over
six policies closes only half of the remaining gap. The ceiling is an estimate, not a bound: the
optimiser could miss a solution, and it only checks staying in the cone, not reaching $\B$. The true
ceiling for $\mathcal{C}_N$ can therefore only be lower.

\begin{table}[H]
\centering\small
\caption{Failure overlap and viability of the shared failures (test split).}
\label{tab:failures}
\tabfit{failures}
\end{table}

\begin{figure}[H]
\centering
\includegraphics[width=\linewidth]{fig_viability.pdf}
\caption{\textbf{Edge test states by fate under \runChosen.} Crashes with no safe trajectory (red) cluster at
small margin and outward speed, at every altitude and vertical speed. The few crashed-but-viable states (amber)
are where the policy could still improve.}
\label{fig:viability}
\end{figure}

\subsection{Actuator commands and smoothness}
\Cref{tab:actuators,fig:actuators} summarise the commands issued while $\pi_{\mathrm{SA}}$ is active
(before the hand-off).
\begin{itemize}
\item \textbf{Saturation} occurs almost only at the first step. The initial thrust command saturates for
\tdThrSatZero\% of test states under TD3 and \sacThrSatZero\% under SAC. Along the trajectory, thrust is at
its maximum for \tdThrMax\% of steps and near zero for \tdThrZero\%, and body rates saturate for \tdRateSat\%.
\item \textbf{Magnitudes.} Commands respect the box by construction. Peak roll/pitch rates reach 16--17\,rad/s at
the 95th percentile, close to $\bar\omega$ but below the SITL rate limit (1100\,deg/s $=19.2$\,rad/s).
\item \textbf{Attitude.} The policy tilts far more than the base controller (\cref{tab:envelope}): median
peak tilt $\tiltMedGen^\circ$ (general) and $\tiltMedEdge^\circ$ (edge), with 95th percentiles of
$\tiltPninetyfiveGen^\circ$ and $\tiltPninetyfiveEdge^\circ$, against $\tiltPninetyfiveGenLQR^\circ$ and
$\tiltPninetyfiveEdgeLQR^\circ$ for $\pi_{\B}$ alone. Large tilts are how it reverses outward velocity
quickly, but they matter for deployment (\cref{sec:limitations}).
\end{itemize}

TD3's actor is steeper than SAC's: its 95th-percentile Jacobian of the normalised action is about 9 vs.\
about 5 (\cref{tab:sens}). What the CIL integrates, however, is the closed-loop sensitivity
$\Phi_t=\partial x_t/\partial x_0$ along the backup rollout, and its constraint rows are
$\nabla\hc(x_t)\Phi_t$. Their peak values over the horizon are similar for all six policies, and the
row magnitudes are slightly smaller for TD3. \texttt{SAC, seed 0} has an outlier state with
$\max_t\|\Phi_t\|=230$. On this evidence the steeper TD3 actor is not a disadvantage for Phase~II.

\begin{table}[H]
\centering\small
\caption{Actuator statistics while $\pi_{\mathrm{SA}}$ is active (all test regions). Saturation means within
2\% of a bound in normalised units. $\|\Delta\tilde u\|$: step-to-step change of the normalised command.}
\label{tab:actuators}
\tabfit{actuators}
\end{table}

\begin{figure}[H]
\centering
\includegraphics[width=\linewidth]{fig_actuators.pdf}
\caption{\textbf{Distribution of commands} while $\pi_{\mathrm{SA}}$ is active, over all test states (log
density). The learned policies use the full thrust range; the base controller stays near hover thrust.}
\label{fig:actuators}
\end{figure}

\begin{table}[H]
\centering\small
\caption{Trajectory envelope of successful arrivals, up to the hand-off.}
\label{tab:envelope}
\tabfit{envelope}
\end{table}

\begin{table}[H]
\centering\small
\caption{Actor smoothness and closed-loop sensitivity. The Jacobian is taken at the 2048 general and edge
test states; $\Phi_t$ at every 8th of them (256 states).}
\label{tab:sens}
\tabfit{sensitivity}
\end{table}

%======================================================================
\section{Closed-loop case studies}\label{sec:cases}
%======================================================================

\Cref{fig:fan} gives an overview: 120 random test states, each rolled out until it enters $\B$ (blue) or
leaves the cone (red). General-region states descend or glide inward to $\B$. Edge states climb and turn
inward within a few tenths of a second. The crashes are the non-viable states on the wall.

\begin{figure}[H]
\centering
\includegraphics[width=\linewidth]{fig_fan.pdf}
\caption{\textbf{Overview of closed-loop behaviour} of the learned backup on 120 random test states,
drawn in the radial plane.}
\label{fig:fan}
\end{figure}

The following cases were selected by rule from the test set to cover distinct situations
(\cref{tab:cases}). Each figure shows, for 3\,s (the 2\,s horizon plus 1\,s of base-controller hold):
\begin{itemize}
\item \textbf{(a)} the path in the radial plane, with the cone and the position projection of $\B$;
\item \textbf{(b)} the top view, with the cone cross-section at the initial altitude;
\item \textbf{(c)} the safety margin $\hc$;
\item \textbf{(d)} the thrust command and \textbf{(e)} the body-rate commands, with their limits;
\item \textbf{(f)} the base-set level $V/c_{\B}$ on a log scale, where entry into $\B$ is the crossing of 1.
\end{itemize}
The grey shading marks the base-controller phase after the hand-off.

\begin{table}[H]
\centering\small
\caption{Case studies: initial condition (altitude, radius, margin, outward and vertical speed, tilt) and
closed-loop outcome under \runChosen{} ($t$: arrival or crash time; minima and maxima over the
safe-arrival segment). Units: m, m/s, degrees, s, m/s$^2$, rad/s.}
\label{tab:cases}
\tabfit{cases}
\end{table}

\subsection*{Case A: high, off-axis start descending toward the wall}
The vehicle starts $\caseAfzeta$\,m above the pad at $\caseAfr$\,m from the axis, $\caseAfh$\,m from the wall,
descending at $\caseAfvzabs$\,m/s. The policy first lets it glide down along the wall: the margin dips to
$\caseAhmin$\,m while it builds inward velocity, with thrust between $\caseAamin$ and 14\,m/s$^2$ and a tilt
of up to $\caseAtilt^\circ$. It then brakes with a thrust pulse of $\caseAamax$\,m/s$^2$ just before
entering $\B$ at $t=\caseAt$\,s, about 1\,m above the pad. After the hand-off, thrust settles at $g$, the
rates at zero, and $V$ decays exponentially: the base set is invariant, as certified.

\begin{figure}[H]
\centering
\includegraphics[width=\linewidth]{fig_case_A.pdf}
\caption{\textbf{Case A} (general region): long descent along the wall, braking pulse, hand-off.}
\label{fig:caseA}
\end{figure}

\subsection*{Case B: low-altitude edge state moving outward and down}
The vehicle starts $\caseBfh$\,m from the wall at $\caseBfzeta$\,m altitude, moving outward at $\caseBfvout$\,m/s
and down at $\caseBfvzabs$\,m/s. The policy commands near-full thrust ($\caseBamax$\,m/s$^2$) for about
0.1\,s while tilting the thrust vector inward (peak tilt $\caseBtilt^\circ$). Because the cone widens with
altitude, climbing increases the margin as effectively as moving inward: after a minimum of $\caseBhmin$\,m
the margin rises steeply. The vehicle enters $\B$ at $t=\caseBt$\,s, and thrust drops to $\caseBamin$\,m/s$^2$
to shed the climb rate just before the hand-off. This climb-and-turn-in manoeuvre is the characteristic
recovery for edge states (\cref{fig:fan}b).

\begin{figure}[H]
\centering
\includegraphics[width=\linewidth]{fig_case_B.pdf}
\caption{\textbf{Case B} (edge band): full-thrust climb and inward roll from 4\,cm off the wall.}
\label{fig:caseB}
\end{figure}

\subsection*{Case C: large tilt, descending close to the pad}
The vehicle starts at $\caseCfzeta$\,m with $\caseCftilt^\circ$ tilt and $\caseCfspeed$\,m/s speed, moving
outward ($\caseCfvout$\,m/s) and down ($\caseCfvzabs$\,m/s). This is the most demanding of the successful
cases. Thrust saturates at $\caseCamax$\,m/s$^2$, the pitch rate reaches $\caseCwmax$\,rad/s, and the vehicle
tilts to $\caseCtilt^\circ$ to stop the outward motion. The margin bottoms out at $\caseChmin$\,m;
arrival is at $t=\caseCt$\,s.

\begin{figure}[H]
\centering
\includegraphics[width=\linewidth]{fig_case_C.pdf}
\caption{\textbf{Case C} (general region): thrust saturation and a hard roll to stop outward motion near the pad.}
\label{fig:caseC}
\end{figure}

\subsection*{Case D: capture shell and base-controller hold}
Starting just outside $\B$ ($V\approx c_{\B}$) at $\caseDfspeed$\,m/s, the safe-arrival policy acts for only
$\caseDt$\,s before the state enters $\B$. The remaining 2.96\,s show the certified base pair at work:
modest commands ($\le\caseDamax$\,m/s$^2$, $\le\caseDwmax$\,rad/s), a margin that never drops below
$\caseDhmin$\,m, and $V$ decreasing by four orders of magnitude.

\begin{figure}[H]
\centering
\includegraphics[width=\linewidth]{fig_case_D.pdf}
\caption{\textbf{Case D} (capture shell): immediate hand-off, then forward invariance and convergence under $\pi_{\B}$.}
\label{fig:caseD}
\end{figure}

\subsection*{Case E: a non-viable state}
The vehicle starts $\caseEfh$\,m from the wall at $\caseEfzeta$\,m altitude, moving outward at $\caseEfvout$\,m/s
and down at $\caseEfvzabs$\,m/s. Descending also shrinks the margin, because the cone narrows downward.
The policy responds correctly and at the limits: thrust $\approx\caseEamax$\,m/s$^2$ (within 3\% of
$a_{\max}$) and roll rate $\caseEwmax$\,rad/s from the first steps. It still leaves the cone at
$t=\caseEt$\,s, by $\caseEhminabs$\,m. The best trajectory-optimisation result (orange) issues almost the same
commands and also exits (by $\caseEopthminabs$\,m). This is the longest-surviving of the \edgeNonviable{}
non-viable edge states; most fail within $\tfMedian$\,s. It illustrates what the ceiling in
\cref{tab:summary} means: no admissible input could have kept this state safe.

\begin{figure}[H]
\centering
\includegraphics[width=\linewidth]{fig_case_E.pdf}
\caption{\textbf{Case E} (edge band, non-viable): the learned policy and the trajectory optimiser both
saturate and both leave the cone within 0.1\,s. The time axis is zoomed.}
\label{fig:caseE}
\end{figure}

\subsection*{Case F: a successful recovery that touches the ground}
The vehicle starts only $\caseFfzeta$\,m above the pad, near the axis, descending at $\caseFfvzabs$\,m/s with
$\caseFftilt^\circ$ tilt. Braking needs at least $v^2/(2\cdot3g)\approx2.3$\,cm even with the thrust vector
already vertical, and the attitude must first be corrected. The path therefore dips to
$\zeta=\caseFzmin$\,m, below the pad plane, before climbing into $\B$ at $t=\caseFt$\,s. The trajectory
never approaches the cone wall ($\min\hc=\caseFhmin$\,m), so in the model this is a safe arrival. In reality
it is a ground contact. This happens in \belowPadGen\% (general) and \belowPadEdge\% (edge) of successful
arrivals; see \cref{sec:limitations}.

\begin{figure}[H]
\centering
\includegraphics[width=\linewidth]{fig_case_F.pdf}
\caption{\textbf{Case F} (general region): recovery that dips 3\,cm below the pad plane; the model has no ground.}
\label{fig:caseF}
\end{figure}

\subsection*{Case G: the same edge state under TD3 and SAC}
From an edge state moving outward at $\caseGfvout$\,m/s, both backbones arrive at the same time
($\caseGt$\,s and $\caseGct$\,s) with similar paths. Their commands differ in detail: TD3 peaks at
$\caseGamax$\,m/s$^2$ and $\caseGwmax$\,rad/s, SAC at $\caseGcamax$\,m/s$^2$ and $\caseGcwmax$\,rad/s.
This is consistent with the aggregate statistics: the two learners converge to closely related policies.

\begin{figure}[H]
\centering
\includegraphics[width=\linewidth]{fig_case_G.pdf}
\caption{\textbf{Case G} (edge band): TD3 (orange) and SAC (blue) from the same initial state. Dashed body
rates in (e) are SAC's.}
\label{fig:caseG}
\end{figure}

\subsection*{Case H: where the learned policy matters}
An edge state $\caseHfh$\,m from the wall at $\caseHfzeta$\,m altitude, moving outward at $\caseHfvout$\,m/s.
The base controller alone (orange) crashes at $t=\caseHct$\,s: in its linear design, lateral braking must
go through attitude, and it commands at most $\caseHcamax$\,m/s$^2$ of thrust. The learned policy commands
$\caseHamax$\,m/s$^2$ immediately and rolls hard (peak tilt $\caseHtilt^\circ$). The large thrust vector
reverses the outward velocity within a few steps, and the vehicle arrives in $\B$ at $t=\caseHt$\,s. States
of this kind account for the gap between the two rows of \cref{fig:slices}.

\begin{figure}[H]
\centering
\includegraphics[width=\linewidth]{fig_case_H.pdf}
\caption{\textbf{Case H} (edge band): the base controller alone leaves the cone; the learned policy recovers.
Dashed body rates in (e) are the base controller's.}
\label{fig:caseH}
\end{figure}

%======================================================================
\section{Limitations and next steps}\label{sec:limitations}
%======================================================================

\begin{enumerate}
\item \textbf{No ground plane.} The cone continues below the pad to its apex, so recoveries may pass below
$\zeta=0$ (\cref{fig:caseF}). The design region also contains states a few centimetres above the pad that
descend at up to 2\,m/s and cannot avoid contact. Remedies: add $h_{\mathrm{ground}}=\zeta\ge0$ as a second
safety row (both in Phase~I and in the CIL), or restrict $\Omega$ near the pad.
\item \textbf{Aggressive attitudes and the rate-loop assumption.} Body rates are assumed to be tracked
instantly, and the policy exploits this with tilts of $60$--$80^\circ$ at the 95th percentile. The SITL
notes already identify rate-loop lag as the main sim-to-SITL gap for the powerloop. Before deployment:
penalise or constrain tilt in Phase~I, lower $\bar\omega$, or add first-order rate dynamics to the model.
\item \textbf{Numerical certificate.} $c_{\mathrm{Lyap}}$ comes from an adversarial search, and invariance
and containment are verified by sampling and simulation (cone containment is exact up to the boundary grid).
The margin $c_{\B}=\cBratio\,c_{\mathrm{Lyap}}$ accounts for this, but it is not a proof. A certified
bound would need a Taylor-remainder argument for $F$ in the chart.
\item \textbf{Cone geometry is a placeholder.} $r_0$ and $\theta$ should come from the landing reference,
which also enables the reference-tube region $R_{\mathrm{tube}}$ of the note.
\path{scripts/certify_landing_base_set.py} recomputes $\zd$ and $c_{\B}$ for new geometry.
\item \textbf{Non-viable mass in $\Omega$.} About 12\% of the edge band cannot be saved by any input. This
costs nothing in the objective, which just gets a constant offset, but it hides progress: recoverability relative
to the viability ceiling (\fracOfCeil\% for the selected policy) is the more informative number.
\item \textbf{Phase~II for landing} is not yet implemented. It needs cone rows and the 9-D base set in
the backup-CBF layer, a landing reference trajectory, and a Phase~II environment.
\end{enumerate}

\paragraph{Recommended checkpoint.} The actor of \runChosen,
\begin{center}\small\path{outputs/landing_phase1_compare/td3_seed0/landing_backup_policy_actor.pkl},\end{center}
has the highest validation score, a test $\hat\mu_w=\muChosen$ and no sensitivity outliers.

%======================================================================
\appendix
\section{Supplementary tables}
%======================================================================

\begin{table}[H]
\centering\small
\caption{LQR gain $K$ ($\delta u=-Ke$). Entries below $5\times10^{-3}$ in magnitude are shown as 0; the
block structure mirrors the four integrator chains.}
\label{tab:gain}
\tabfit{gain}
\end{table}

\begin{table}[H]
\centering\small
\caption{Level bounds versus hover height ($z$-independent bounds: $c_{\U}=\cU$, $c_{\mathrm{chart}}=\cChart$,
$c_{\mathrm{Lyap}}=\cLyap$).}
\label{tab:zsweep}
\input{tables/zsweep}
\end{table}

\section{Reproduction}
All numbers, tables and figures in this report are generated from the training outputs by
\begin{verbatim}
python docker_batch/report_landing_phase1.py \
    --root outputs/landing_phase1_compare --run td3_seed0
\end{verbatim}
which writes \texttt{figures/}, \texttt{tables/}, \texttt{numbers.tex} and this document, and compiles it
when \texttt{latexmk} is available. The prose was written for the runs reported here.

\begin{thebibliography}{9}
\bibitem{ps2rl} K.~S.~Yun, Z.~Li and N.~Azizan. \emph{Provably Safe, yet Scalable Reinforcement Learning.}
arXiv:2606.14536, 2026.
\bibitem{landingnote} \emph{Reformulating PS2-RL Phase I for the Approach-Cone Landing Task.} Internal note,
28 September 2026.
\end{thebibliography}

\end{document}
